
\documentclass[11pt]{article}

\usepackage{amsmath,amssymb,amsfonts,bm}
\usepackage[a4paper,margin=1in]{geometry}
\usepackage{mathtools}
\usepackage{hyperref}
\usepackage{physics}
\setlength{\emergencystretch}{4em}

% Chebyshev
\newcommand{\T}{\mathcal{T}}
\newcommand{\U}{\mathcal{U}}

\title{Exact propagators and first-passage generating functions on a 1D ring\\
with one directed long-range link: \\
\large (A) lazy-reservoir shortcut \quad vs \quad (B) nearest-neighbour rewiring shortcut\\
\large Giuggioli--Bristol notation: $\underline{B}$ (ring), $\underline{A}$ (defect), $\underline{X}=\underline{B}-\underline{A}$, defect parameters $\eta$}
\author{}
\date{}

\begin{document}
\sloppy
\maketitle

\section*{Overview}
We derive, in full detail:
\begin{itemize}
\item the defect-free generating function (Green's function) $\tilde Q_{n_0}(n,z)$ for a \emph{lazy} nearest-neighbour random walk on a finite ring of size $N$;
\item the exact defective generating function $\tilde S_{n_0}(n,z)$ after introducing a \emph{single directed} long-range link using two constructions:
\begin{enumerate}
\item[(A)] \textbf{Lazy-reservoir shortcut}: the shortcut probability is taken from the self-loop at the source node;
\item[(B)] \textbf{Nearest-neighbour rewiring}: the shortcut probability is taken from an existing nearest-neighbour jump at the source node.
\end{enumerate}
\item two complementary derivations for $\tilde S$: (i) a low-rank (Sherman--Morrison) resolvent update; (ii) the Giuggioli defect-technique determinant formula (scalar for (A), explicit $2\times2$ for (B));
\item the first-passage generating function $\tilde F_{n_0\to n}(z)=\tilde S_{n_0}(n,z)/\tilde S_n(n,z)$, including a fully explicit Chebyshev closed form for construction (A);
\item a detailed comparison of the two constructions and their relation to Watts--Strogatz small-world rewiring.
\item numerical validation against direct matrix inversion (double precision).
\end{itemize}

Throughout we use the Giuggioli--Bristol defect-technique notation ($\underline{B},\underline{A},\underline{X},\eta$), and we keep \emph{all} intermediate steps, including the explicit $k=0$ cancellation in the Chebyshev identity step.

\section{Notation and model}

\subsection{Ring, indexing, and transition matrices}

We consider a discrete-time random walk on a 1D ring (cycle) of size $N$ with periodic boundary conditions (PBC).
Sites are labeled $n\in\{0,1,\dots,N-1\}$ and all indices are understood modulo $N$.

We use the column-stochastic convention
\[
\begin{aligned}
\Pr(X_{t+1}=n\mid X_t=n') &= B_{n,n'} &&\text{(defect-free ring)},\\
\Pr(X_{t+1}=n\mid X_t=n') &= A_{n,n'} &&\text{(with defect)}.
\end{aligned}
\]
Thus $P(t+1)=\underline{B}\,P(t)$ for the defect-free process, and $\sum_n B_{n,n'}=\sum_n A_{n,n'}=1$ for each column $n'$.

The defect matrix is
\[
\underline{X}\equiv \underline{B}-\underline{A}.
\]
In the Supplement notation, each defect \emph{pair} $(u_i,v_i)$ carries two parameters $\eta_{u_i,v_i}$ and $\eta_{v_i,u_i}$, which correspond to the off-diagonal entries of $\underline{X}$:
\[
\eta_{u,v}\equiv X_{v,u}=B_{v,u}-A_{v,u},\qquad \eta_{v,u}\equiv X_{u,v}=B_{u,v}-A_{u,v}.
\]

\subsection{Baseline (defect-free) lazy nearest-neighbour ring}
Let $q\in(0,1)$ be the \emph{total} nearest-neighbour jump probability (nearest-neighbour range $k=1$).
Then the defect-free transition matrix $\underline{B}$ is
\[
B_{n,n}=1-q,\qquad B_{n\pm 1,n}=\frac{q}{2},\qquad B_{m,n}=0\ \text{otherwise}.
\]
In words: at each step, with probability $q$ the walker moves to one of the two neighbours (left/right with probability $q/2$ each),
and with probability $1-q$ the walker stays in place (self-loop / sojourn).

\subsection{Propagators and generating functions (Green's functions)}
\begin{itemize}
\item Defect-free propagator: $Q_{n_0}(n,t)=\Pr(X_t=n\mid X_0=n_0)$ under $\underline{B}$.
\item Defective propagator: $S_{n_0}(n,t)=\Pr(X_t=n\mid X_0=n_0)$ under $\underline{A}$.
\item Generating function (discrete-time $z$-transform): $\tilde f(z)=\sum_{t=0}^{\infty} f(t)\,z^t$.
\end{itemize}
We write $\tilde Q_{n_0}(n,z)=\sum_{t\ge 0} Q_{n_0}(n,t)\,z^t$ and $\tilde S_{n_0}(n,z)$ similarly.

Matrix-wise,
\[
\underline{G}_B(z)\equiv (\underline{I}-z\underline{B})^{-1}=\sum_{t=0}^\infty z^t \underline{B}^t,
\qquad
\underline{G}_A(z)\equiv (\underline{I}-z\underline{A})^{-1}=\sum_{t=0}^\infty z^t \underline{A}^t,
\]
and therefore
\[
\tilde Q_{n_0}(n,z)=[\underline{G}_B(z)]_{n,n_0},\qquad
\tilde S_{n_0}(n,z)=[\underline{G}_A(z)]_{n,n_0}.
\]

\subsection{Ring distance and Chebyshev polynomials}
For any two sites $a,b$ define the minimal ring distance
\[
d(a,b)\equiv \min\bigl(|a-b|,\;N-|a-b|\bigr)\in\Bigl\{0,1,\dots,\Bigl\lfloor \frac N2\Bigr\rfloor\Bigr\}.
\]
Chebyshev polynomials of the first and second kind are
\[
\T_n(x)=\cos(n\arccos x),\qquad
\U_n(x)=\frac{\sin((n+1)\arccos x)}{\sin(\arccos x)}.
\]

\section{Part I --- Defect-free propagator $\tilde Q_{n_0}(n,z)$ in Chebyshev form}

\subsection{Step 1: $z$-transform of the master equation (fully explicit)}

The defect-free master equation is
\begin{equation}
Q_{n_0}(n,t+1)=\frac q2\,Q_{n_0}(n-1,t)+\frac q2\,Q_{n_0}(n+1,t)+(1-q)\,Q_{n_0}(n,t),
\label{eq:master_ring}
\end{equation}
with initial condition $Q_{n_0}(n,0)=\delta_{n,n_0}$ and PBC.

Multiply both sides of \eqref{eq:master_ring} by $z^t$ and sum $t=0$ to $\infty$:
\begin{align*}
\sum_{t=0}^{\infty} Q_{n_0}(n,t+1) z^t
&=\frac q2\sum_{t=0}^{\infty} Q_{n_0}(n-1,t) z^t
 +\frac q2\sum_{t=0}^{\infty} Q_{n_0}(n+1,t) z^t
 +(1-q)\sum_{t=0}^{\infty} Q_{n_0}(n,t) z^t.
\end{align*}
The sums on the right are the generating functions:
\[
\sum_{t=0}^{\infty} Q_{n_0}(n\pm 1,t) z^t=\tilde Q_{n_0}(n\pm 1,z),
\qquad
\sum_{t=0}^{\infty} Q_{n_0}(n,t) z^t=\tilde Q_{n_0}(n,z).
\]
For the left-hand side we use the time-shift identity:
\begin{align*}
\sum_{t=0}^{\infty} Q_{n_0}(n,t+1) z^t
&=\frac{1}{z}\sum_{t=0}^{\infty} Q_{n_0}(n,t+1) z^{t+1}
=\frac{1}{z}\sum_{s=1}^{\infty} Q_{n_0}(n,s) z^{s}\\
&=\frac{1}{z}\left(\sum_{s=0}^{\infty} Q_{n_0}(n,s) z^{s}-Q_{n_0}(n,0)\right)
=\frac{1}{z}\left(\tilde Q_{n_0}(n,z)-\delta_{n,n_0}\right).
\end{align*}
Substituting into the summed master equation yields
\[
\frac{1}{z}\left(\tilde Q_{n_0}(n,z)-\delta_{n,n_0}\right)
=
\frac q2\,\tilde Q_{n_0}(n-1,z)+\frac q2\,\tilde Q_{n_0}(n+1,z)
+(1-q)\,\tilde Q_{n_0}(n,z).
\]
Multiplying by $z$ and rearranging gives the Helmholtz-type equation
\begin{equation}
\bigl[1-z(1-q)\bigr]\tilde Q_{n_0}(n,z)-\frac{qz}{2}\Bigl(\tilde Q_{n_0}(n-1,z)+\tilde Q_{n_0}(n+1,z)\Bigr)
=\delta_{n,n_0}.
\label{eq:helmholtz_ring}
\end{equation}

\subsection{Step 2: discrete Fourier representation on the ring}

Let
\[
\theta_k\equiv \frac{2\pi k}{N},\qquad k=0,1,\dots,N-1.
\]
The eigenvalues of the ring transition matrix $\underline{B}$ are
\[
\lambda_k=1-q+q\cos\theta_k.
\]
Therefore the generating function (Green's function) is
\begin{equation}
\tilde Q_{n_0}(n,z)
=\frac{1}{N}\sum_{k=0}^{N-1}\frac{e^{i\theta_k(n-n_0)}}{1-z\lambda_k}
=
\frac{1}{N}\sum_{k=0}^{N-1}\frac{e^{i\theta_k(n-n_0)}}{1-z\,[1-q+q\cos\theta_k]}.
\label{eq:Q_fourier}
\end{equation}
By translational invariance, $\tilde Q_{n_0}(n,z)$ depends only on $d=d(n,n_0)$.

\subsection{Step 3: rewrite the sum in the form required by Identity E3}

Introduce the spectral parameter
\begin{equation}
\sigma(z)\equiv \frac{1-z(1-q)}{qz}.
\label{eq:sigma_def}
\end{equation}
Then
\[
1-z[1-q+q\cos\theta_k]=qz\bigl(\sigma(z)-\cos\theta_k\bigr).
\]
Insert this into \eqref{eq:Q_fourier} and use the ring symmetry (pair $k$ with $N-k$) to write the numerator as a cosine:
\begin{equation}
\tilde Q_{n_0}(n,z)
=
\frac{1}{qz}\,\frac{1}{N}\sum_{k=0}^{N-1}\frac{\cos\!\bigl(d\,\theta_k\bigr)}{\sigma(z)-\cos\theta_k},
\qquad d=d(n,n_0).
\label{eq:Q_cos_sum}
\end{equation}

\subsection{Step 4: apply Identity E3 and show the $k=0$ cancellation explicitly}

Identity E3 (periodic case, Appendix~E in the 2020 reference) states that for $m\in\{0,1,\dots,N\}$,
\begin{equation}
\sum_{k=1}^{N-1}\frac{\cos\!\left(\frac{2\pi m k}{N}\right)}{\sigma-\cos\!\left(\frac{2\pi k}{N}\right)}
=
\frac{1}{1-\sigma}
+
N\,\frac{\T_{N-m}(\sigma)+\T_m(\sigma)}{(\sigma^2-1)\,\U_{N-1}(\sigma)}.
\label{eq:E3_identity}
\end{equation}

\paragraph{Align the cosine.}
Using $\theta_k=\frac{2\pi k}{N}$, we write
\[
\cos(d\theta_k)=\cos\!\left(d\frac{2\pi k}{N}\right)=\cos\!\left(\frac{2\pi(m=d)\,k}{N}\right).
\]
Thus we apply \eqref{eq:E3_identity} with
\[
m=d,\qquad \sigma=\sigma(z),
\]
which yields for the partial sum $k=1,\dots,N-1$:
\begin{equation}
\sum_{k=1}^{N-1}\frac{\cos(d\theta_k)}{\sigma(z)-\cos\theta_k}
=
\frac{1}{1-\sigma(z)}
+
N\,\frac{\T_{N-d}(\sigma(z))+\T_{d}(\sigma(z))}{(\sigma(z)^2-1)\,\U_{N-1}(\sigma(z))}.
\label{eq:partial_sum}
\end{equation}

\paragraph{Restore the $k=0$ term.}
We now write the full sum from \eqref{eq:Q_cos_sum} as
\[
\sum_{k=0}^{N-1}\frac{\cos(d\theta_k)}{\sigma(z)-\cos\theta_k}
=
\underbrace{\frac{\cos(0)}{\sigma(z)-\cos(0)}}_{\text{$k=0$ term}}
+
\sum_{k=1}^{N-1}\frac{\cos(d\theta_k)}{\sigma(z)-\cos\theta_k}.
\]
Since $\cos(0)=1$, the missing $k=0$ term equals $1/(\sigma(z)-1)$. Using \eqref{eq:partial_sum},
\begin{align}
\sum_{k=0}^{N-1}\frac{\cos(d\theta_k)}{\sigma(z)-\cos\theta_k}
&=
\frac{1}{\sigma(z)-1}
+
\frac{1}{1-\sigma(z)}
+
N\,\frac{\T_{N-d}(\sigma(z))+\T_{d}(\sigma(z))}{(\sigma(z)^2-1)\,\U_{N-1}(\sigma(z))}.
\label{eq:full_sum_with_cancel}
\end{align}
The first two terms cancel explicitly:
\[
\frac{1}{\sigma(z)-1}+\frac{1}{1-\sigma(z)}
=
\frac{1}{\sigma(z)-1}-\frac{1}{\sigma(z)-1}=0.
\]
Therefore the full sum becomes
\begin{equation}
\sum_{k=0}^{N-1}\frac{\cos(d\theta_k)}{\sigma(z)-\cos\theta_k}
=
N\,\frac{\T_{N-d}(\sigma(z))+\T_{d}(\sigma(z))}{(\sigma(z)^2-1)\,\U_{N-1}(\sigma(z))}.
\label{eq:full_sum_final}
\end{equation}

Substitute \eqref{eq:full_sum_final} back into \eqref{eq:Q_cos_sum} to obtain the closed form:
\begin{equation}
\boxed{
\tilde Q_{n_0}(n,z)=
\frac{\T_{d}(\sigma(z))+\T_{N-d}(\sigma(z))}
{qz\,\bigl(\sigma(z)^2-1\bigr)\,\U_{N-1}(\sigma(z))},
\qquad d=d(n,n_0).
}
\label{eq:Q_chebyshev}
\end{equation}

\subsection{Odd vs even $N$ (explicit)}
Let $d=d(n,n_0)\in\{0,1,\dots,\lfloor N/2\rfloor\}$.
\begin{itemize}
\item If $N$ is odd, \eqref{eq:Q_chebyshev} holds as written for all $d$.
\item If $N$ is even and $d=N/2$ (antipodes), then $N-d=d$ and the numerator becomes $2\T_{N/2}(\sigma(z))$:
\[
\tilde Q_{n_0}(n,z)\big|_{d=N/2}
=
\frac{2\,\T_{N/2}(\sigma(z))}
{qz\,\bigl(\sigma(z)^2-1\bigr)\,\U_{N-1}(\sigma(z))}.
\]
\end{itemize}

\section{Part II --- One directed shortcut: two constructions and their exact $\tilde S$}

\subsection{A unifying algebraic view: ``probability-mass transfer'' in one column}

Both constructions we consider are \emph{local} modifications of a single column of $\underline{B}$:
we pick a \emph{source} node $u$, and move probability mass $p$ from some \emph{old destination} $b$ to a \emph{new destination} $a$.
At the level of transition matrices this is
\begin{equation}
\boxed{
\underline{A}=\underline{B}+p\,(\ket{a}-\ket{b})\bra{u}.
}
\label{eq:A_rank1}
\end{equation}
Here $\ket{a}$ denotes the basis vector at site $a$.

\paragraph{Validity constraints.}
The only constraint is that the modified transition probabilities remain nonnegative.
If $b$ is the self-loop ($b=u$), then $p\le B_{u,u}=1-q$.
If $b$ is a nearest neighbour ($b=u\pm1$), then $p\le B_{u\pm1,u}=q/2$.

\subsection{Sherman--Morrison derivation of the resolvent update (all steps)}

Define the defect-free and defective resolvents
\[
\underline{G}_B(z)=(\underline{I}-z\underline{B})^{-1},\qquad \underline{G}_A(z)=(\underline{I}-z\underline{A})^{-1}.
\]
Using \eqref{eq:A_rank1}, we write
\[
\underline{I}-z\underline{A}=\underline{I}-z\underline{B}-z p(\ket a-\ket b)\bra u
\equiv \underline{M}-\bm{u}\bm{v}^\top,
\]
with
\[
\underline{M}\equiv \underline{I}-z\underline{B},
\qquad
\bm{u}\equiv z p(\ket a-\ket b),
\qquad
\bm{v}^\top \equiv \bra u.
\]
We now invert $\underline{M}-\bm{u}\bm{v}^\top$ exactly.

\paragraph{Claim (Sherman--Morrison for a rank-1 update).}
If $\underline{M}$ is invertible and $1-\bm{v}^\top \underline{M}^{-1}\bm{u}\neq 0$, then
\begin{equation}
(\underline{M}-\bm{u}\bm{v}^\top)^{-1}
=
\underline{M}^{-1}
+
\frac{\underline{M}^{-1}\bm{u}\bm{v}^\top \underline{M}^{-1}}{1-\bm{v}^\top \underline{M}^{-1}\bm{u}}.
\label{eq:SM}
\end{equation}

\paragraph{Verification.}
Let $\underline{R}$ denote the right-hand side of \eqref{eq:SM}.
Compute $(\underline{M}-\bm{u}\bm{v}^\top)\underline{R}$:
\begin{align*}
(\underline{M}-\bm{u}\bm{v}^\top)\underline{R}
&=(\underline{M}-\bm{u}\bm{v}^\top)\underline{M}^{-1}
+(\underline{M}-\bm{u}\bm{v}^\top)\frac{\underline{M}^{-1}\bm{u}\bm{v}^\top \underline{M}^{-1}}{1-\bm{v}^\top \underline{M}^{-1}\bm{u}}\\
&=\underline{I}-\bm{u}\bm{v}^\top \underline{M}^{-1}
+\frac{\bm{u}\bm{v}^\top \underline{M}^{-1}}{1-\bm{v}^\top \underline{M}^{-1}\bm{u}}
-\frac{\bm{u}\bm{v}^\top \underline{M}^{-1}\bm{u}\bm{v}^\top \underline{M}^{-1}}{1-\bm{v}^\top \underline{M}^{-1}\bm{u}}.
\end{align*}
Factor $\bm{u}\bm{v}^\top \underline{M}^{-1}$ in the last two terms:
\[
\frac{\bm{u}\bm{v}^\top \underline{M}^{-1}}{1-\bm{v}^\top \underline{M}^{-1}\bm{u}}
\left[1-(\bm{v}^\top \underline{M}^{-1}\bm{u})\right]
=
\bm{u}\bm{v}^\top \underline{M}^{-1}.
\]
Thus the last two terms together equal $+\bm{u}\bm{v}^\top \underline{M}^{-1}$, which cancels the $-\bm{u}\bm{v}^\top \underline{M}^{-1}$ term.
Therefore $(\underline{M}-\bm{u}\bm{v}^\top)\underline{R}=\underline{I}$, proving \eqref{eq:SM}.

\paragraph{Apply to the random walk.}
Since $\underline{M}^{-1}=(\underline{I}-z\underline{B})^{-1}=\underline{G}_B(z)$, \eqref{eq:SM} gives
\begin{equation}
\underline{G}_A(z)
=
\underline{G}_B(z)
+
\frac{z p\,\underline{G}_B(z)\,(\ket a-\ket b)\bra u\,\underline{G}_B(z)}
{1-z p\,\bra u\,\underline{G}_B(z)\,(\ket a-\ket b)}.
\label{eq:GA_update}
\end{equation}
Take the matrix element $(n,n_0)$ to obtain the defective propagator generating function:
\begin{equation}
\boxed{
\tilde S_{n_0}(n,z)
=
\tilde Q_{n_0}(n,z)
+
\frac{z p\,[\tilde Q_a(n,z)-\tilde Q_b(n,z)]\,\tilde Q_{n_0}(u,z)}
{1-z p\,[\tilde Q_a(u,z)-\tilde Q_b(u,z)]}.
}
\label{eq:S_general_ab}
\end{equation}
This formula is \emph{exact} and will be specialised to the two constructions below.

\subsection{Construction (A): shortcut probability taken from the self-loop (lazy reservoir)}

\paragraph{Definition of the defect.}
We add a directed shortcut $u\to v$ of probability $p$ by taking $p$ from the self-loop at $u$:
\[
A_{v,u}=B_{v,u}+p=p,\qquad
A_{u,u}=B_{u,u}-p=(1-q)-p,
\qquad
A_{u\pm1,u}=B_{u\pm1,u}=\frac q2.
\]
This is \eqref{eq:A_rank1} with
\[
a=v,\qquad b=u,
\qquad 0\le p\le 1-q.
\]

\paragraph{Exact propagator (rank-1 form).}
Specialising \eqref{eq:S_general_ab} to $(a=v,b=u)$ gives
\begin{equation}
\boxed{
\tilde S^{\text{(A)}}_{n_0}(n,z)
=
\tilde Q_{n_0}(n,z)
+
\frac{z p\,[\tilde Q_v(n,z)-\tilde Q_u(n,z)]\,\tilde Q_{n_0}(u,z)}
{1-z p\,[\tilde Q_v(u,z)-\tilde Q_u(u,z)]}.
}
\label{eq:S_A_uv}
\end{equation}
Using translational invariance on the ring, $\tilde Q_x(y,z)=\tilde Q(d(x,y),z)$, this can be written entirely in terms of distances:
\begin{equation}
\boxed{
\tilde S^{\text{(A)}}_{n_0}(n,z)
=
\tilde Q(d(n,n_0),z)
-
\frac{z p\,[\tilde Q(d(n,u),z)-\tilde Q(d(n,v),z)]\,\tilde Q(d(u,n_0),z)}
{1+z p\,[\tilde Q(0,z)-\tilde Q(d(u,v),z)]}.
}
\label{eq:S_A_dist}
\end{equation}

\paragraph{Defect-technique encoding (one defect pair).}
In terms of $\underline{X}=\underline{B}-\underline{A}$, the only nonzero off-diagonal entry is
\[
X_{v,u}=-p,
\]
so the Supplement defect parameter is
\[
\eta_{u,v}=X_{v,u}=-p,\qquad \eta_{v,u}=0,
\]
and the diagonal change $X_{u,u}=+p$ is implied by column-stochasticity.

\paragraph{Recovering (A) from the defect-technique determinant formula ($M=1$).}
Although (A) is already solved by the rank-1 update \eqref{eq:S_A_uv}--\eqref{eq:S_A_dist}, it is useful to connect to the
Giuggioli defect-technique formula (especially since (B) will naturally appear as ``multiple defects'').
For $M$ defect pairs $(u_i,v_i)$, the defect-technique result is
\begin{equation}
\tilde S_{n_0}(n,z)
=
\tilde Q_{n_0}(n,z)-1+\frac{\det\!\bigl[\underline{H}(n,n_0;z)\bigr]}{\det\!\bigl[\underline{H}(z)\bigr]}.
\label{eq:defect_formula}
\end{equation}
For a single defect pair $(u,v)$ we have $M=1$ and the determinants reduce to scalars.
With the standard definitions
\begin{align}
\det[\underline{H}(z)]
&=
H(z)
=
\eta_{u,v}\,\tilde Q_{(u-v)}(u,z)-\eta_{v,u}\,\tilde Q_{(u-v)}(v,z)-\frac{1}{z},
\label{eq:H_scalar}\\[1mm]
\det[\underline{H}(n,n_0;z)]
&=
H(n,n_0;z)
=
H(z)-\tilde Q_{(u-v)}(n,z)\Bigl[\eta_{u,v}\,\tilde Q_{n_0}(u,z)-\eta_{v,u}\,\tilde Q_{n_0}(v,z)\Bigr],
\label{eq:Hn_scalar}
\end{align}
where $\tilde Q_{(u-v)}(x,z)\equiv \tilde Q_u(x,z)-\tilde Q_v(x,z)$,
and inserting $\eta_{u,v}=-p$, $\eta_{v,u}=0$ gives
\begin{align}
\det[\underline{H}(z)]
&=
-\frac{1}{z}
-p\Bigl(\tilde Q(0,z)-\tilde Q(d(u,v),z)\Bigr),
\label{eq:H_A}\\[2mm]
\det[\underline{H}(n,n_0;z)]
&=
\det[\underline{H}(z)]
+p\Bigl(\tilde Q(d(n,u),z)-\tilde Q(d(n,v),z)\Bigr)\,\tilde Q(d(u,n_0),z).
\label{eq:Hn_A}
\end{align}
Substituting \eqref{eq:H_A}--\eqref{eq:Hn_A} into \eqref{eq:defect_formula} reproduces exactly \eqref{eq:S_A_dist}.

\subsection{Construction (B): rewiring probability from a nearest-neighbour jump to the shortcut}

\paragraph{Definition of the defect.}
Choose one nearest neighbour of $u$ as the \emph{old destination} $w$; for definiteness on the ring we can take $w=u+1$
(the ``clockwise'' neighbour). We create a directed shortcut $u\to v$ of probability $p$ by \emph{reducing} the probability of the jump $u\to w$ by $p$:
\[
\begin{aligned}
A_{v,u} &= B_{v,u}+p=p,\\
A_{w,u} &= B_{w,u}-p=\frac q2-p,\\
A_{u,u} &= B_{u,u}=1-q,\\
A_{u-1,u} &= B_{u-1,u}=\frac q2.
\end{aligned}
\]
This is \eqref{eq:A_rank1} with
\[
a=v,\qquad b=w=u+1,\qquad 0\le p\le \frac q2.
\]
(The case $p=q/2$ corresponds to a \emph{full} rewiring of the entire clockwise edge probability onto the shortcut.)

\paragraph{Exact propagator (rank-1 form).}
Specialising \eqref{eq:S_general_ab} to $(a=v,b=w)$ gives
\begin{equation}
\boxed{
\tilde S^{\text{(B)}}_{n_0}(n,z)
=
\tilde Q_{n_0}(n,z)
+
\frac{z p\,[\tilde Q_v(n,z)-\tilde Q_w(n,z)]\,\tilde Q_{n_0}(u,z)}
{1-z p\,[\tilde Q_v(u,z)-\tilde Q_w(u,z)]}.
}
\label{eq:S_B_uvw}
\end{equation}
In distance form (using $d(u,w)=1$ on the ring):
\begin{equation}
\boxed{
\tilde S^{\text{(B)}}_{n_0}(n,z)
=
\tilde Q(d(n,n_0),z)
+
\frac{z p\,[\tilde Q(d(n,v),z)-\tilde Q(d(n,w),z)]\,\tilde Q(d(u,n_0),z)}
{1-z p\,[\tilde Q(d(u,v),z)-\tilde Q(1,z)]}.
}
\label{eq:S_B_dist}
\end{equation}

\subsection{How construction (B) appears as ``multiple defects'' in the $\eta$-pair formalism (explicit $2\times2$ determinant)}

Although \eqref{eq:S_B_uvw} is rank-1 at the matrix level, in the Supplement pair notation it is convenient to represent (B) as the superposition of \emph{two} defect pairs, because the self-loop at $u$ is \emph{unchanged} and must emerge as a cancellation.

\paragraph{Two defect pairs.}
Write (B) as ``add $p$ to $u\to v$ and subtract $p$ from the self-loop'' (pair 1) \emph{plus}
``subtract $p$ from $u\to w$ and add $p$ to the self-loop'' (pair 2).
This corresponds to two pairs:
\[
(u_1,v_1)=(u,v),\qquad (u_2,v_2)=(u,w),
\]
with defect parameters
\[
\eta_{u_1,v_1}=\eta_{u,v}=-p,\qquad \eta_{u_2,v_2}=\eta_{u,w}=+p,
\qquad
\eta_{v_1,u_1}=\eta_{v,u}=0,\qquad \eta_{v_2,u_2}=\eta_{w,u}=0.
\]
Indeed, this yields off-diagonal changes
\[
X_{v,u}=\eta_{u,v}=-p,\qquad X_{w,u}=\eta_{u,w}=+p,
\]
and diagonal change $X_{u,u}=-(\eta_{u,v}+\eta_{u,w})=0$ (so $A_{u,u}=B_{u,u}$), as required.

\paragraph{Build $\underline{H}(z)$ explicitly.}
For $M=2$ defect pairs, the Supplement defines $\underline{H}(z)$ by
\[
\begin{aligned}
H(z)_{i,j}
&=
\eta_{u_i,v_i}\,\tilde Q_{(u_j-v_j)}(u_i,z)
-\eta_{v_i,u_i}\,\tilde Q_{(u_j-v_j)}(v_i,z)
-\frac{\delta_{i,j}}{z},\\
\tilde Q_{(u_j-v_j)}(x,z)
&=\tilde Q_{u_j}(x,z)-\tilde Q_{v_j}(x,z).
\end{aligned}
\]
Here $\eta_{v_i,u_i}=0$, and $u_1=u_2=u$, so
\[
\tilde Q_{(u-v)}(u,z)=\tilde Q_u(u,z)-\tilde Q_v(u,z)=\tilde Q(0,z)-\tilde Q(d(u,v),z),
\]
\[
\tilde Q_{(u-w)}(u,z)=\tilde Q_u(u,z)-\tilde Q_w(u,z)=\tilde Q(0,z)-\tilde Q(d(u,w),z)=\tilde Q(0,z)-\tilde Q(1,z).
\]
Therefore
\[
\underline{H}(z)=
\begin{pmatrix}
-\dfrac{1}{z}-p\bigl(\tilde Q(0,z)-\tilde Q(d(u,v),z)\bigr)
&
-p\bigl(\tilde Q(0,z)-\tilde Q(1,z)\bigr)
\\[2mm]
+p\bigl(\tilde Q(0,z)-\tilde Q(d(u,v),z)\bigr)
&
-\dfrac{1}{z}+p\bigl(\tilde Q(0,z)-\tilde Q(1,z)\bigr)
\end{pmatrix}.
\]
Its determinant is computed step-by-step:
\begin{align*}
\det[\underline{H}(z)]
&=\frac{1}{z^2}
+\frac{p}{z}\Bigl[(\tilde Q(0)-\tilde Q_{uv})-(\tilde Q(0)-\tilde Q_{1})\Bigr]\\
&=\frac{1}{z^2}
+\frac{p}{z}\bigl(\tilde Q_{1}-\tilde Q_{uv}\bigr).
\end{align*}
where $\tilde Q_{uv}\equiv \tilde Q(d(u,v),z)$ and $\tilde Q_{1}\equiv \tilde Q(1,z)$.
The $p^2$ terms cancel exactly.
Thus
\begin{equation}
\boxed{
\det[\underline{H}(z)]
=
\frac{1}{z^2}\Bigl[1+z p\bigl(\tilde Q(1,z)-\tilde Q(d(u,v),z)\bigr)\Bigr]
=
\frac{1}{z^2}\Bigl[1-z p\bigl(\tilde Q_{uv}-\tilde Q_{1}\bigr)\Bigr].
}
\label{eq:detH_B}
\end{equation}

\paragraph{Build $\underline{H}(n,n_0;z)$ and its determinant.}
The Supplement defines
\[
H(n,n_0;z)_{i,j}
=
H(z)_{i,j}
-\tilde Q_{(u_j-v_j)}(n,z)\Bigl[\eta_{u_i,v_i}\tilde Q_{n_0}(u_i,z)-\eta_{v_i,u_i}\tilde Q_{n_0}(v_i,z)\Bigr].
\]
Since $\eta_{v_i,u_i}=0$ and $u_i=u$ for both $i$, the bracket reduces to $\eta_{u_i,v_i}\tilde Q_{n_0}(u,z)$.
Define
\[
\begin{aligned}
Q_{u n_0}(z) &\equiv \tilde Q_{n_0}(u,z),\\
\Delta_v(n) &\equiv \tilde Q_{(u-v)}(n,z)=\tilde Q_u(n,z)-\tilde Q_v(n,z),\\
\Delta_w(n) &\equiv \tilde Q_{(u-w)}(n,z)=\tilde Q_u(n,z)-\tilde Q_w(n,z).
\end{aligned}
\]
Then
\[
\underline{H}(n,n_0;z)=
\begin{pmatrix}
H_{11}+p\,\Delta_v(n)\,Q_{u n_0}(z) & H_{12}+p\,\Delta_w(n)\,Q_{u n_0}(z)\\
H_{21}-p\,\Delta_v(n)\,Q_{u n_0}(z) & H_{22}-p\,\Delta_w(n)\,Q_{u n_0}(z)
\end{pmatrix},
\]
and a direct expansion gives (keeping every term)
\begin{align*}
\det[\underline{H}(n,n_0;z)]
&=\det[\underline{H}(z)]
+\frac{p}{z}\,Q_{u n_0}(z)\,\bigl(\Delta_w(n)-\Delta_v(n)\bigr).
\end{align*}
Finally note that $\Delta_w(n)-\Delta_v(n)=\tilde Q_v(n,z)-\tilde Q_w(n,z)$, so
\begin{equation}
\boxed{
\det[\underline{H}(n,n_0;z)]
=
\frac{1}{z^2}\Bigl[1-z p\bigl(\tilde Q_{uv}-\tilde Q_{1}\bigr)\Bigr]
+\frac{p}{z}\,\tilde Q_{n_0}(u,z)\,\bigl(\tilde Q_v(n,z)-\tilde Q_w(n,z)\bigr).
}
\label{eq:detHn_B}
\end{equation}

\paragraph{Recover $\tilde S^{\mathrm{(B)}}$ from the determinant formula.}
The defect technique gives
\[
\tilde S_{n_0}(n,z)=\tilde Q_{n_0}(n,z)-1+\frac{\det[\underline{H}(n,n_0;z)]}{\det[\underline{H}(z)]}.
\]
Substituting \eqref{eq:detH_B}--\eqref{eq:detHn_B} yields
\[
-1+\frac{\det[\underline{H}(n,n_0;z)]}{\det[\underline{H}(z)]}
=
\frac{z p\,\tilde Q_{n_0}(u,z)\,(\tilde Q_v(n,z)-\tilde Q_w(n,z))}{1-z p(\tilde Q_{uv}-\tilde Q_{1})},
\]
which reproduces exactly \eqref{eq:S_B_uvw}.

\section{Chebyshev packaging for the defective propagators}

Define
\[
\sigma=\sigma(z),\qquad
\mathcal{C}(z)\equiv qz\bigl(\sigma^2-1\bigr)\U_{N-1}(\sigma),\qquad
\mathcal{N}_d(\sigma)\equiv \T_d(\sigma)+\T_{N-d}(\sigma),
\]
so that \eqref{eq:Q_chebyshev} is
\[
\tilde Q(d,z)=\frac{\mathcal{N}_d(\sigma)}{\mathcal{C}(z)}.
\]
Then every term in \eqref{eq:S_A_dist} and \eqref{eq:S_B_dist} is explicit in $\T_n,\U_{N-1}$.

For example, in construction (B) the denominator becomes
\[
1-z p\bigl(\tilde Q(d(u,v),z)-\tilde Q(1,z)\bigr)
=
1-\frac{z p}{\mathcal{C}(z)}\Bigl(\mathcal{N}_{d(u,v)}(\sigma)-\mathcal{N}_{1}(\sigma)\Bigr),
\]
and similarly for the numerator differences.

\section{Part III --- First-passage generating function}

For either construction, the first-passage generating function to site $n$ starting from $n_0$ is
\begin{equation}
\boxed{
\tilde F_{n_0\to n}(z)=\frac{\tilde S_{n_0}(n,z)}{\tilde S_{n}(n,z)}.
}
\label{eq:F_def}
\end{equation}

\subsection{Construction (A): explicit form}
Insert \eqref{eq:S_A_dist} for the numerator and the same formula with $n_0=n$ for the denominator:
\[
\tilde F^{\text{(A)}}_{n_0\to n}(z)=\frac{\tilde S^{\text{(A)}}_{n_0}(n,z)}{\tilde S^{\text{(A)}}_{n}(n,z)}.
\]
All occurrences of $\tilde Q$ are given explicitly by \eqref{eq:Q_chebyshev}.

\paragraph{Chebyshev closed form (construction (A)).}
Starting from \eqref{eq:F_def} with \eqref{eq:S_A_dist}, define the two recurring blocks
\begin{equation}
D_Q(z)\equiv 1+z p\bigl(\tilde Q(0,z)-\tilde Q(d(u,v),z)\bigr),
\qquad
\Delta_n(z)\equiv \tilde Q(d(n,u),z)-\tilde Q(d(n,v),z).
\label{eq:DQ_Delta_def}
\end{equation}
Multiplying numerator and denominator by $D_Q(z)$ yields the equivalent form without nested fractions:
\begin{equation}
\boxed{
\tilde F^{\text{(A)}}_{n_0\to n}(z)=
\frac{
\tilde Q(d(n,n_0),z)\,D_Q(z)-z p\,\Delta_n(z)\,\tilde Q(d(u,n_0),z)
}{
\tilde Q(0,z)\,D_Q(z)-z p\,\Delta_n(z)\,\tilde Q(d(u,n),z)
}.
}
\label{eq:F_A_no_nested}
\end{equation}
Now substitute the Chebyshev packaging from \eqref{eq:Q_chebyshev}:
\[
\tilde Q(d,z)=\frac{\mathcal{N}_d(\sigma(z))}{\mathcal{C}(z)},
\qquad
\mathcal{C}(z)\equiv qz\bigl(\sigma(z)^2-1\bigr)\U_{N-1}(\sigma(z)),
\qquad
\mathcal{N}_d(\sigma)\equiv \T_d(\sigma)+\T_{N-d}(\sigma),
\]
and introduce the compact factor
\begin{equation}
\boxed{
D(z)\equiv \mathcal{C}(z)+z p\bigl(\mathcal{N}_0(\sigma(z))-\mathcal{N}_{d(u,v)}(\sigma(z))\bigr),
}
\label{eq:D_def}
\end{equation}
so that $D_Q(z)=D(z)/\mathcal{C}(z)$. Cancelling the common $\mathcal{C}(z)^{-2}$ factor gives the fully explicit closed form
\begin{equation}
\boxed{
\tilde F^{\text{(A)}}_{n_0\to n}(z)=
\frac{
\mathcal{N}_{d(n,n_0)}(\sigma(z))\,D(z)
-z p\,\Bigl(\mathcal{N}_{d(n,u)}(\sigma(z))-\mathcal{N}_{d(n,v)}(\sigma(z))\Bigr)\,\mathcal{N}_{d(u,n_0)}(\sigma(z))
}{
\mathcal{N}_{0}(\sigma(z))\,D(z)
-z p\,\Bigl(\mathcal{N}_{d(n,u)}(\sigma(z))-\mathcal{N}_{d(n,v)}(\sigma(z))\Bigr)\,\mathcal{N}_{d(u,n)}(\sigma(z))
}.
}
\label{eq:F_A_cheb_closed}
\end{equation}
Here $\mathcal{N}_0(\sigma)=\T_0(\sigma)+\T_N(\sigma)=1+\T_N(\sigma)$.

\subsection{Construction (B): explicit form}
Insert \eqref{eq:S_B_dist} for the numerator and the same formula with $n_0=n$ for the denominator:
\[
\tilde F^{\text{(B)}}_{n_0\to n}(z)=\frac{\tilde S^{\text{(B)}}_{n_0}(n,z)}{\tilde S^{\text{(B)}}_{n}(n,z)}.
\]
Again, all terms reduce to Chebyshev polynomials via \eqref{eq:Q_chebyshev}.

\section{Part IV --- Detailed comparison and relation to Watts--Strogatz SWN rewiring}

\subsection{Mechanistic comparison: where does the shortcut probability come from?}

\paragraph{(A) Lazy-reservoir shortcut (self-loop transfer).}
\begin{itemize}
\item Probability mass is taken from the self-loop at the source: $u\to u$ decreases, $u\to v$ increases.
\item Constraint: $0\le p\le 1-q$.
\item Nearest-neighbour probabilities $u\to u\pm1$ remain exactly $q/2$, so the local diffusive step statistics are unchanged.
\item In the $\eta$-pair formalism this is one defect pair: $M=1$.
\end{itemize}

\paragraph{(B) Nearest-neighbour rewiring shortcut (edge-to-edge transfer).}
\begin{itemize}
\item Probability mass is taken from an existing jump $u\to w$ (with $w=u+1$ or $u-1$) and moved to $u\to v$.
\item Constraint: $0\le p\le q/2$ for partial rewiring of a single neighbour edge; the full rewiring limit is $p=q/2$.
\item The self-loop probability at $u$ remains $1-q$, but the local jump kernel at $u$ becomes biased if only one neighbour is rewired.
\item In the $\eta$-pair formalism it is naturally represented by \emph{two} defect pairs $(u,v)$ and $(u,w)$ with opposite signs, so $M=2$,
even though the net matrix update is still rank-1.
\end{itemize}

\subsection{Relation to Watts--Strogatz SWN}

In a classical Watts--Strogatz small-world network (SWN), one rewires edges of an underlying ring lattice.
If the random walk chooses uniformly among neighbours, then the transition probabilities are proportional to $1/\chi_n$ (node degree).
A single rewiring event changes the degree of (at least) the target node and the removed node, hence \emph{all} outgoing transition probabilities
from those nodes must be renormalised, producing \emph{multiple} nonzero entries in $X=B-A$ and thus multiple defects in the determinant formalism.
In contrast, constructions (A) and (B) above operate at the level of \emph{transition probabilities} directly:
they move a prescribed probability mass $p$ within one column, so they remain low-rank and analytically tractable.

\subsection{Consistency checks}

\begin{itemize}
\item $p=0$ in either construction gives $\underline{A}=\underline{B}$ and $\tilde S=\tilde Q$.
\item In (A), the maximal shortcut is $p=1-q$ (all sojourn probability at $u$ is converted into shortcut jumps).
\item In (B), the full rewiring of the clockwise edge corresponds to $p=q/2$ (the entire $u\to w$ probability is redirected to $u\to v$).
\end{itemize}

\section*{Numerical validation (double precision)}

We validated \eqref{eq:Q_chebyshev}, \eqref{eq:S_A_dist}, and \eqref{eq:F_A_cheb_closed} against direct matrix computation of the resolvent
for construction (A) (lazy-reservoir shortcut).
For each parameter set $(N,q,p,u,v,z)$ we built the column-stochastic transition matrices $\underline{B}$ (ring) and $\underline{A}$ (directed shortcut),
and computed
\[
G_B(z)=(\underline{I}-z\underline{B})^{-1},\qquad
G_A(z)=(\underline{I}-z\underline{A})^{-1}.
\]
The ground-truth generating functions are $\tilde Q_{n_0}(n,z)=[G_B(z)]_{n,n_0}$ and $\tilde S_{n_0}(n,z)=[G_A(z)]_{n,n_0}$.
For first passage we used $\tilde F_{n_0\to n}(z)=[G_A(z)]_{n,n_0}/[G_A(z)]_{n,n}$ for $n_0\neq n$.
We report the maximum absolute error over all indices:
\[
\max_{n,n_0}\bigl|\Delta\tilde Q\bigr|,\qquad
\max_{n,n_0}\bigl|\Delta\tilde S\bigr|,\qquad
\max_{n_0\neq n}\bigl|\Delta\tilde F\bigr|.
\]

\begin{center}
\small
\begin{tabular}{r r r r r r l l l}
\hline
$N$ & $q$ & $p$ & $u$ & $v$ & $z$ & $\max|\Delta\tilde Q|$ & $\max|\Delta\tilde S|$ & $\max|\Delta\tilde F|$ \\
\hline
9 & 0.6 & 0.2 & 1 & 4 & 0.3 & 2.22e-16 & 4.44e-16 & 5.55e-17 \\
10 & 0.7 & 0.1 & 2 & 7 & 0.5 & 4.44e-16 & 2.22e-16 & 5.55e-17 \\
12 & 0.5 & 0.3 & 0 & 6 & 0.6 & 4.44e-16 & 4.44e-16 & 8.33e-17 \\
15 & 0.8 & 0.05 & 3 & 4 & 0.4 & 2.22e-16 & 4.44e-16 & 1.11e-16 \\
20 & 0.4 & 0.2 & 5 & 13 & 0.75 & 1.33e-15 & 1.33e-15 & 1.67e-16 \\
\hline
\end{tabular}
\end{center}

In a sweep of 200 random cases ($N\le 40$, seed $=0$, $q\in[0.2,0.9]$, $z\in[0.2,0.85]$, $p\in[0,1-q]$), we found
\[
\max|\Delta\tilde Q|=3.11\times 10^{-15},\qquad
\max|\Delta\tilde S|=3.11\times 10^{-15},\qquad
\max|\Delta\tilde F|=4.44\times 10^{-16}.
\]
These residuals are consistent with floating-point round-off and the conditioning of matrix inversion as $z\to 1^{-}$.

\end{document}
